% commentaire.tex — f172v, vv. 838–875 (conversion initiale des travaux antérieurs)
\comm{838}{σκάφιον ἀποκεκαρμένην.}{C'est l'humiliation que le tyran de Cumes, Aristodème, avait fait subir aux filles de son pays : περιτρόχαλα κείρεσθαι, tandis qu'il forçait les garçons à porter des cheveux longs. Plutarque, \italique{De mul. virtut}, 26, p. 361f. (Note de Willems.)\\
La scholie précise à propos de σκάφιον : \citComm{εἶδος κουρᾶς δουλικῆς.\\
\francais{Genre de coupe rase qui convient à l'esclave.}}}
\comm{839}{ὦ πόλεις.}{Dans son livre \italique{Aristophanes und Eupolis} (2007), Natalia Kyriakidi a montré qu'Eupolis, le poète comique rival d'Aristophane, est selon toute vraisemblance mort en 411 lors de la victoire navale de Κυνὸς Σῆμα (cf. comm. vv. 833-7). Pour arriver à cette conclusion, elle récuse très justement la plupart des témoignages, qui exhalent l'anecdote recyclée, pour ne conserver que le témoignage de Suidas (ε 3657) -- ou de la \italique{Souda}, si l'on préfère :\citComm{Εὔπολις, Σωσίπολις, Ἀθηναῖος, κωμικός. ὃς ἑπτακαίδεκα ἔτη γεγονὼς ἤρξατο ἐπιδείκνυσθαι: καὶ ἐδίδαξε μὲν δράματα ιζ, ἐνίκησε δὲ ζ. καὶ ἀπέθανε ναυαγήσας κατὰ τὸν Ἑλλήσποντον ἐν τῷ πρὸς Λακεδαιμονίους πολέμῳ: καὶ ἐκ τούτου ἐκωλύθη στρατεύεσθαι ποιητήν. \\
Fils de Sosipolis, Athénien, un poète comique. Il a commencé à monter des pièces de théâtre à l'âge de dix-sept ans. Il a produit 17 pièces et en a gagné 7. Il est mort dans le naufrage de son navire dans l'Hellespont pendant la guerre contre les Lacédémoniens ; à la suite de cela, il fut interdit à un poète de servir en campagne. }Natalia continue (p. 11) : \allemand{Die zwei großen Schlachten, die am Hellespont während des Peloponnesischen Krieges stattfanden, sind der Sieg der Athener bei} Κυνὸς Σῆμα \allemand{411 und die endgültige Niederlage bei }Αἰγὸς Ποταμοί \allemand{405. Der Tod des Eupolis während einer Seeschlacht könnte auch die Entstehung der Geschichte vom Ertränken des Eupolis durch Alkibiades erklären} [$=$ Eupolis, selon certaines légendes, aurait été noyé en mer par Alcibiade au retour de Sicile], \allemand{weil beide Geschichten als Vorfälle im Meer bei einer Kriegshandlung wenige Jahre vor dem Ende des Peloponnesischen Krieges überliefert sind. Storey vermutet, dass Ailians aitiologische Erklärung des Ortsnamens} Κυνὸς Θρῆνος [$=$ le lieu commun du chien qui meurt de chagrin sur la tombe de son maître] \allemand{auf }Κυνὸς Σῆμα \allemand{zurückzuführen ist. Da kein Werk des Eupolis nach 411 zu datieren ist und für die Zeit danach keine Indizien für eine weitere Tätigkeit des Dichters überliefert sind, erscheint die Information der Suda als das einzige zuverlässige Zeugnis über den Tod des Eupolis. So enden sein Leben und seine dichterische Karriere wahrscheinlich im Sommer 411 während der Seeschlacht bei} Κυνὸς Σῆμα.\\
Pour le texte, voir le commentaire critique.
S'il faut attribuer l'idée d'accorder les sièges des Athéniennes avec la valeur de leur progéniture à feu Eupolis (et peut-être aussi celle de les tondre le cas décevant échéant), il n'aura pas manqué de s'indigner lui aussi que la mère d'Hyperbolos spécule avec le fruit de la rapine de son fils. 
On ne possède malheureusement que peu de choses d'Eupolis. On sait que le \italique{Maricas} dut être représenté en $-421$, que le chœur était formé de deux demi-chœurs : celui des riches et celui des pauvres citoyens, chœurs qui débattaient avec le démagogue Hyperbolos.
La comédie est une violente charge contre ce dernier.
Elle se terminait par une danse licencieuse de la mère du démagogue.
Nombre de commentateurs, à l'instar des scholiastes, ont rapproché les \italique{Cavaliers} du \italique{Maricas}... mais pas les \italique{Thesmophories} -- plus pauvres en scholies anciennes.
Pourtant, certaines idées d'Eupolis semblent en être reprises. Nous l'allons voir tout-à-l'heure.
Mettons en garde tout-de-suite : les hypothèses -- comme les explications -- seront fort spéculatives.\\
Un commentaire du \italique{Maricas}, daté du \nieme{II} siècle ap. J.-C., nous est parvenu de façon très morcelaire. On y lit cette citation (fr. 192.96-99 Kassel-Austin) :\citComm{τοῦτ᾽ ἐκδανείζει καὶ κυκᾷς [τὸ]\douteux{ν} ναυτικὸν \douteux{ἄ}[ργυρον \ \ \ \ \ (\textbf{ἄργυρον} Bibi)\\
ἐπὶ τῷ πέμπτῳ μέρε[ι \ \ \ \ \ \ \ ]τοὺς τόκ[ους οἱ\\
ναυτικοί, ἡμεῖς δ᾽ ἄρ᾽ οἴκ[ \ \ \ \ \ \ \ \ \ \ \ ]εν· ὁ τῶν\\
πλουσίων λέγει χορός.}Le chœur des riches (ὁ τῶν πλουσίων χορός) reproche à Maricas ($=$ Hyperbolos) de prêter (ἐκδανείζει) en concoctant (κυκᾷς) des prêts (τὸν ναυτικὸν ἄργυρον) au cinquième (ἐπὶ τῷ πέμπτῳ μέρει\nc{Sens bien attesté d'ἐπί. Cf. Dém. 1250, 20 : ἐπ᾽ ὀκτὼ ὀβολοῖς δανείζειν τοῦ μηνὸς τὴν μνᾶν, «~prêter de l'argent moyennant un intérêt de huit oboles la mine par mois~», c'est-à-dire à $16 \%$.}), c'est-à-dire avec un taux d'intérêt de $20 \%$.
Il est question ensuite des τόκους, `enfants' ($=$ intérêts) de ces ναυτικοί (cf. vers 845).
Continuant le jeu de mot nautique (les prêts sont littéralement de l'argent \italique{marin}), est-ce à ce moment qu'Eupolis, par la voix du chœur, proposait de faire tondre \italique{nautiquement} cet usurier ou sa mère et de lui faire porter une \italique{barque} (σκάφιον vers 838) ?\\
Revenons aux \italique{Thesmophories} : en guise d'oraison funèbre, Aristophane rend hommage à la pertinence des propositions d'Eupolis :\citComm{τῷ γὰρ εἰκὸς\nc{Le tour elliptique apparaît aussi au vers 703 des \italique{Acharniens} où il est glosé `πῶς δίκαιόν ἐστιν' par une scholie. Le ton est effectivement indigné : «~comment peut-il être juste/raisonnable que...~», c'est-à-dire «~quelle indignité que...~». Voir aussi Sophocle (\italique{Ajax}, vers 1366-1367), non indigné mais ironique :\citComm{Ἀγαμέμνων -- ἡ πάνθ᾽ ὅμοια πᾶς ἀνὴρ αὑτῷ πονεῖ.\\
Ὀδυσσεύς -- τῷ γάρ με μᾶλλον εἰκὸς ἢ ᾽μαυτῷ πονεῖν;\\
Agamemnon -- En toute chose, chacun ne travaille que pour lui-même.\\
Ulysse -- Comment peut-il être plus raisonnable que de travailler pour moi-même ?}} ωὖπολι\\
Comment tolérer, ô Eupolis, que \&c.}Dernier coup de chapeau d'Aristophane à son défunt rival, mort des suites de la folie des Athéniens et de l'incompétence des amiraux.}
\comm{840–1}{Ὑπερβόλου ... Λαμάχου,}{Hyperbolos, le démagogue athénien -- marchand de lampes de son état\nc{\italique{Cavaliers}, v. 739 (non nommé) ; \italique{Nuées}, v. 1064 ; \italique{Paix}, v. 690.} -- devint démagogue en chef après la mort de Cléon devant Amphipolis en 422 av. J.-C.\nc{Voir la \italique{Paix}, v. 679-692, Paix qui fout le camp, révoltée que le peuple se soit donné un si indigne patron.} ; sa mère et lui furent lors une cible privilégiée d'Aristophane et des poètes comiques\nc{Cf. les \italique{Nuées}, v. 551-562 : Aristophane daube sur ses plagiaires de contemporains, notamment Eupolis réchauffant leurs \italique{Cavaliers} pour son \italique{Maricas} et y rajoutant un autre plat déjà concocté par Phrynicos en son temps : une mère d'Hyperbolos, ivrognesse et dansant le licencieux cordax. Sur les attaques des comiques contre la mère d'Hyperbolos, présentée comme une débauchée notoire, voir aussi les fragments 7-12 des \italique{Vendeuses de pain} (αἱ Ἀρτοπώλιδες) d'Hermippos (dans l'édition Kassel-Austin) ainsi que les fragments de Platon le Comique dans sa pièce éponyme \italique{Hyperbolos}.} (suite sans doute à quelque lapsus à relent freudien du démagogue, lapsus qui fut repris à l'encan par les Comiques). Cet ὀξίνην Ὑπέρβολον\nc{\italique{Cavaliers}, v. 1303. Thucydide, qui n'estimait guère plus ce μοχθηρὸν ἄνθρωπον, ne parle qu'une seule fois d'Hyperbolos : pour nous annoncer sa mort et en brosser un portrait expéditif (VIII, 73.3) : \citComm{καὶ Ὑπέρβολόν τέ τινα τῶν Ἀθηναίων, μοχθηρὸν ἄνθρωπον, ὠστρακισμένον οὐ διὰ δυνάμεως καὶ ἀξιώματος φόβον, ἀλλὰ διὰ πονηρίαν καὶ αἰσχύνην τῆς πόλεως, ἀποκτείνουσι μετὰ Χαρμίνου\\
S'en prenant à Hyperbolos, un Athénien méprisable qu'on avait ostracisé non par peur de son influence et de son prestige, mais parce que c'était un malhonnête homme qui déshonorait la cité, ils [les conjurés favorables à l'oligarchie] l'assassinèrent avec l'aide de Charminos.}} réunissait tout pour être détesté d'Aristophane : de caractère exécrable, il était un partisan, comme son mentor Cléon, de la guerre à outrance\nc{\italique{Cavaliers}, v. 1303-1315 ; il cherchait à renforcer la flotte pour que l'expédition de Sicile fût poussée jusqu'à Carthage.} ; c'était un intenteur de procès\nc{\italique{Acharniens}, v. 846-847.}, un fourbe\nc{\italique{Guêpes}, v. 1007.}, qui s'est enrichi par ses coquineries\nc{\italique{Nuées}, v. 1064-1066. Tel fils, telle mère ? L'usure brocardée par Aristophane dont fait preuve la mère d'Hyperbolos n'était-elle pas la continuation des activités louches du fils, depuis que ce dernier se trouvait en exil à Samos ?}. Il eut l'insigne honneur d'être le dernier des Athéniens ostracisés en 417\nc{Plut., \italique{Nicias} 11 ; \italique{Alcib.} 13.}. Il fut assassiné à Samos, son lieu d'exil\nc{Thucydide, VIII, 73.3.}, et son corps jeté en mer\nc{Théopompe 115F69, in \italique{Die Fragmente der griechischen Historiker}, Teil 2, \italique{Zeitgeschichte - B}.}... cette \italique{inhudation} lui valut une dernière raillerie de la part d'Aristophane\nc{\italique{Grenouilles}, v. 570.}. 
\textbf{Lamachos} était un stratège athénien\nc{Thucydide IV, 75. 1 : il commande une dizaine de navires dans le Pont.} et l'un des représentants athéniens de la Paix de Nicias\nc{Thucydide V 19.2 et 24.1}. Avec Nicias et Alcibiade, il est investi des pleins pouvoirs lors de l'expédition de Sicile\nc{Thucydide (VI, 8) ; il s'opposa à Alcibiade sur la stratégie (V 49) : Lamachos voulait aller droit au but et attaquer Syracuse sitôt le débarquement fait ; il finit par se ranger à l'opinion d'Alcibiade (VI, 50).} où il trouve la mort\nc{Thucydide (VI, 101) : Lamachos courant à la rescousse de l'aile droite de l'armée athénienne se retrouva isolé avec quelques-uns de ses camarades ; il fut tué par les Syracusains qui récupérèrent son corps avant de le restituer (VI, 103) aux Athéniens. Nicias fut dès lors le seul commandant en Sicile.}. Lamachos présente une originalité unique parmi les personnages historiques d'Aristophane puisque le poète opère un volte-face à son endroit : du va-t-en-guerre au panache ridicule dépeint dans les \italique{Acharniens} et la \italique{Paix}\nc{Lamachos arborait sur son casque trois aigrettes qui lui donnaient des faux airs de Gorgone (\italique{Acharniens}, v. 574 ; v. 964-965). Aristophane rapproche cet accoutrement d'épouvante à d'autre monstre-croquemitaine au nom voisin (Μορμώ, \italique{Acharniens} v. 582 ; \italique{Paix} v. 474). Enfin, cette crinière en bataille -- de circonstance pour un militaire ! -- permettait un jeu analogue en grec entre λόφοι et λόχοι, trop tentant pour qu'Aristophane n'y cédât pas (\italique{Acharniens}, v. 567, 574, 1074).} (accusé même de s'enrichir grâce à la guerre\nc{Il est traité de μισθαρχίδης au v. 597 des \italique{Acharniens} ; attaque renouvelée au vers 619.}), Aristophane lui rend hommage ici comme dans les \italique{Grenouilles}\nc{\italique{Grenouilles}, v. 1039.}. La mort héroïque de Lamachos\nc{Voir -- outre le passage de Thucydide (VI, 101) -- Plutarque, \italique{Nicias} 18 (et aussi le \italique{Lachès} (19c) de Platon où il est question du courage de Lamachos.)} est sans doute pour beaucoup dans ce changement radical d'opinion.}
\comm{845}{ἀξία γοῦν εἶ τόκου τεκοῦσα τοιοῦτον τόκον.}{Jeu de mots qui roule sur le double sens de τόκος, \latin{usura} et \latin{proles}. C'est ainsi que l'orateur Eschine disait de Démosthène qui, souffleté par Midias, s'était désisté moyennant 30 mines : \citComm{ὁ γὰρ ἄνθρωπος οὐ κεφαλήν, ἀλλὰ πρόσοδον κέκτηται.\\
Ce n'est pas une tête (κεφαλήν), mais un capital que possède cet homme.}(Cf. \italique{Contre Ctésiphon}, 212 ; note de Willems.)}
\comm{850}{τὴν καινὴν Ἑλένην.}{« La nouvelle \italique{Hélène} ». \italique{Nouvelle} parce qu'Euripide l'avait présentée il y a peu de temps, sous l'archontat de Cléocrite (printemps 412). Mais l'adjectif καινή est aussi à prendre au sens d'\italique{innovante}, \italique{imprévue}, \italique{étrange}. Voici ce que dit Victor-Henry Debidour de cette « \italique{Hélène} nouveau style »\nc{Dans l'introduction à sa traduction de l'\italique{Hélène}. Oui ! car Debidour le grand traducteur d'Aristophane, qui a su rendre un peu de la verve de notre poète dans la langue de Rabelais, Debidour ne s'est pas limité à la Comédie Ancienne ; il est aussi l'auteur d'une traduction des pièces conservées des trois grands Tragiques : \italique{Les Tragiques -- Eschyle, Sophocle, Euripide, théâtre complet}, La Pochothèque (1999). Traduction à laquelle s'ajoute l'excellent dossier et les index fort utiles de Paul Demont et Anne Lebeau. Victor-Henry Debidour vérifie ainsi le paradoxe de Socrate dont Platon n'a pas cru utile de nous donner la démonstration à la fin de son \italique{Banquet}.} :\\
Dans la tradition légendaire, Hélène de Sparte, femme de Ménélas, était le parangon de la beauté. Mais aussi le plus scandaleux exemple d'infidélité féminine, cause, par son inconduite, des malheurs sans nombre de la guerre de Troie. Le personnage s'était stéréotypé en « femme fatale », d'autant plus haïssable qu'elle était plus prestigieuse ; et \italique{Euripide lui-même ne la traite pas autrement} -- sauf justement dans la tragédie très peu tragique qui porte son nom. Certes, elle y reste, bien-sûr, la belle Hélène, et si charmeuse que le fils du sage Protée l'aime comme un fou, et comme un sot. Mais la voici transformée en épouse sans tache, victime innocente de la longue perfidie des deux déesses, dont l'une l'a prostituée et l'autre séquestrée : elles l'ont engagée dans une rocambolesque histoire où elle mérite toute notre sympathie. Pour se faire désigner par Pâris comme « la plus belle », Cypris avait promis au prince-berger troyen de lui donner Hélène ; Héra, furieuse, contre sa rivale et son juge, fit en sorte que Pâris n'emmenât à Troie qu'un fantôme d'Hélène, une « nuée » toute semblable à la vraie et suffisamment consistante pour leurrer le ravisseur pendant les dix ans qu'il la garda, et le légitime époux après ce qu'il crut être sa reconquête. La vraie Hélène, elle, miraculeusement transportée en Égypte, y était tenue cachée durant toutes ces années, confiée à la vertueuse protection de Protée. Ainsi, c'est pour un « crime sans criminelle » que des milliers de preux se sont entre-tués et que Troie a été anéantie.\\
Telle est l'\italique{étonnante} variante que recueille Euripide sur la plus célèbre entreprise guerrière de la légende grecque. Car il ne l'a pas inventée : Hérodote en fait état bien avant lui\nc{Hérodote II, 112-120. Par la suite, du reste, le thème de la réhabilitation d'Hélène fut maintes fois repris en exercice d'école : la paradoxe devient à son tour lieu commun.}. De ce point de départ, Euripide a tiré un canevas sur lequel il a brodé une aventure romanesque et humoristique qui souvent ferait croire qu'avant d'être parodiée elle-même par Aristophane, l'\italique{Hélène} donnait dans la parodie.}
\comm{852}{τί κοικύλλεις ἔχων;}{Κοικύλλειν est un verbe à redoublement expressif\nc{Sur le même modèle que μοιμύλλειν, δενδίλλειν, δανδαίνειν, παπταίνειν, \&c.} qui tire son origine de τὰ κύλα, \italique{les poches sous les yeux}, d'où le sens d'« écarquiller les yeux »\nc{Cf. Hésychius κ 3232 : κοικύλλειν· περιβλέπειν· κύλα γὰρ τὰ ἐπάνω τῶν ὀφθαλμῶν [$=$ la partie supérieure des paupières, mais voir aussi la glose d'Hésychius sous κύλα, où le terme désigne la partie \italique{inférieure} des paupières]. οἱ δὲ τὸ περιστρέφεσθαι καὶ διατρίβειν περιβλεπόμενον. ἄλλοι φθονεῖν.}. Voici pour le sens en lequel l'emploie Critylla.
Le verbe est particulièrement bien adapté à la situation : notre anti-héros vient de lancer `à la mer' ses `rames' dans l'espoir de rameuter Euripide ; il est là maintenant, le regard hagard, devant cette mer imaginaire, à attendre... pareil à ce personnage comique, surnommé Κοικυλίων, type de l'idiot qui comptait les vagues de la mer\nc{Ælien, \italique{Varia Historia}, XIII, 15 ($=$ \latin{Adespota comica}, fr. 72 K.-A.) : 
\citComm{φασὶ παχύτατον γενέσθαι τὴν διάνοιαν οἱ τῆς κωμῳδίας ποιηταὶ Πολύδωρόν τινα, καὶ ἄλλον Κοικυλίωνα ὄνομα, ὅσπερ τὰ κύματα ἠρίθμει ὑπὸ τῆς ἄγαν ἀνοίας.\\
\francais{Les poètes comiques disent qu'un certain Polydore (aussi [connu sous] un autre nom : Cœcylion) affolit tant qu'il comptait les vagues sous le coup de son trop plein de folie.}}Ce passage, déniché par Fritzsche (à l'érudition et la bibliothèque toujours impressionnantes !), est à rapprocher du chapitre du Περὶ βλασφημιῶν de Suétone (édité, pour ce qu'il a pu en reconstituer, en 1967 aux Belles Lettres par le grand Taillardat) consacré à quelques idiots célèbres. Περὶ βλασφημιῶν, β 189 $=$ Eustathe \italique{Comm. Od.} p. 1669.53 : ὥσπερ δὲ τούτους ἡ ἱστορία λόγου ἠξίωσεν ἐν οἷς καὶ τὸν τὰ κύματα μετροῦντα Πολύ[δ]ωρον.}, le comble de l'inutilité pour un grec : κύματα μετρεῖν était même passé en proverbe (détourné par Aristophane au vers 1257 de la \italique{Paix}).}
\comm{853}{πικρὰν ἑλένη.}{Le détournement du \italique{Palamède} ayant échoué (après celui du \italique{Télèphe}), Mnésiloque songe maintenant à recourir à la récente \italique{Hélène} d'Euripide. Seulement, de toutes les \italique{Thesmophories}, aucun des geôliers de Mnésiloque ne rentre jamais dans le jeu de ses parodies d'Euripide. Pour le texte, voir le commentaire critique.
Kritylla menace Mnésiloque d'un bon coup de flambeau, flambeau qu'elle avait lors de la procession et l'arrivée au temple. (Coup de flambeau qu'elle assénera au vers 858 en évoquant Hécate \textbf{φωσ}φόρος.)}
\comm{855–8}{Début du prologue de l'\italique{Hélène}.}{Ainsi commence la véritable pièce d'Euripide :
\citComm{Ἑλένη --- Νείλου μὲν αἵδε καλλιπάρθενοι ῥοαί,\\
ὃς ἀντὶ δίας ψακάδος Αἰγύπτου πέδον\\
λευκῆς \difference{τακείσης χιόνος ὑγραίνει γύας.}\\
\francais{Du Nil voici les rives aux belles Nymphes,\\
Nil qui à la place de la pluie dispensée par Zeus\\
arrose la plaine d'Égypte une fois la blanche neige fondue et irrigue les terres.}}Comme tous les semi-habiles – euphémisme pascalien : comprenez `semi-débiles' – Euripide aime à bavasser ; il ne peut s'empêcher dès le troisième vers de la pièce d'exposer son opinion sur la raison des crues du Nil – Ah ça ! se croit-il obligé d'en avoir une quel que soit le sujet ? et croit-il par-dessus le marché qu'elle nous intéresse ? Selon donc le très savant Euripide, ces crues sont dues à la fonte des neiges. Théorie fantaisiste et traitée avec dédain par Hérodote (II, 19-25), condescendant vis-à-vis de cette thèse soutenue par Anaxagore\nc{59A91 D.-K. = fr. 300 Eschyle (ed. Radt). Eschyle avait en son temps aussi évoqué cette fonte des neiges dans les \italique{Suppliantes}, v. 559 sq.}, le maître d'Euripide – raison pour laquelle ce dernier se croit tenu de la défendre ? 
En tout cas il récidive, imperturbable, quelques années après l'\italique{Hélène} et les \italique{Thesmophories}, l'an 408/407, dans \italique{Archélaos}\nc{Fr. 1 \italique{Ἀρχέλαος} Jouan-van Looy = fr. 228 ed. Kannicht.}, tragédie aujourd'hui perdue. (Remarquons que dans le prologue d'\italique{Hélène} comme dans le fragment d'\italique{Archélaos}, les ῥοαί ne doivent pas se traduire par \italique{flots} mais bien par \italique{rives} ou \italique{berges}\nc{Hésychius (ρ 415) : ῥοία· κυλίστρα τῶν ἵππων παρὰ τῷ ποταμῷ καὶ ψάμμῳ. Rien que pour l'\italique{Hélène}, les ῥοαί sont à prendre au sens de rives au moins aux vers 53, 124 (avec jeu de mots sur Εὐρώτας), 250, 492.} et que l'expression « καλλιπάρθενοι ῥοαί » suggère peut-être un jeu de mots avec les Ῥοαί, les nymphes des grenadiers\nc{Eustathe, \italique{Commentaire à l'Odyssée}, p. 1572, 35 : ῥοιαὶ ὁμωνύμως τῷ δένδρῳ καὶ αἱ περὶ αὐτὰς μυθικαὶ νύμφαι.}.)

Attaquons-nous au détournement du prologue d'\italique{Hélène} par Aristophane (\italique{Thesm.} vv. 855-857). Pour traîner dans la boue – ou pis – la muse d'Euripide, notre poète ne recule devant rien. La syrméa dont il est question au vers 857 était une espèce de rave très abondante en Égypte. On en tirait un suc portant le même nom utilisé comme purgatif\nc{Souda (σ 1667) : Συρμαία : τὸ καθαρτικὸν δι' ἐμέτων καὶ κοιλίας. καὶ Συρμαίζειν τὸ κενοῦσθαι, ἐκ ζειῶν καὶ ὕδατος. οὕτως Ἡρόδοτος. λάχανον σελίνῳ ἐοικός : ἢ βρῶμα διὰ στέατος καὶ μέλιτος.\\
Et voici ce qu'en dit Hésychius (σ 2782) : σύρμαια· […] ἔστι δὲ βρωμάτιον διὰ στέατος καὶ μέλιτος. λέγεται δὲ καὶ συρμαϊσμός, καὶ ἔστι πρὸς κάθαρσιν. καὶ πόμα δι' ὕδατος καὶ ἁλῶν.}. Suivant Hérodote\nc{Hérodote II, 77 ; et aussi II, 88 ; II, 125.}, les Égyptiens se purgeaient menstruellement, trois jours de suite. Ce curieux hygiénisme est assez plaisant pour qu'Aristophane n'ait pas manqué de le railler.

Aristophane avait certainement mis en scène une confrontation entre le vendeur de syrméa et le vendeur de pâtisseries dans ses \italique{Danaïdes} (comédie aujourd'hui perdue) – les gâteaux de miel étaient en effet réputés astringents\nc{C'est ce qui ressort d'un passage des \italique{Nuées} (vers 506-508) : Strepsiade hésite à entrer dans l'antre de Socrate, qu'il compare aimablement à celui de Trophonius : εἰς τὼ χεῖρέ νυν / δός μοι μελιτοῦτταν πρότερον, ὡς δέδοικ' ἐγὼ / εἴσω καταβαίνων ὥσπερ εἰς Τροφωνίου. « Mets-moi d'abord en main un gâteau au miel ; car j'ai peur de descendre, comme si je pénétrais dans l'antre de Trophonios. » Il faut bien sûr prononcer d'un trait `εἰς Τροφωνίου' pour faire entendre στρόφος, la colique (la peur se manifestant chez les Grecs par des troubles du bas-ventre ; en français : \italique{avoir les pétoches}). En passant, Euripide suggère le même jeu de mots scatologique : au vers 300 de son \italique{Ion} : Créuse dit à Ion qu'elle est venue avec son mari, mais qu'il a dû s'arrêter en chemin (ἐκστρέφει Τροφωνίου ; Euripide, ainsi qu'il est coutumier, joue à différents niveaux de compréhension)… Le tour était proverbial et Érasme l'a recensé (\italique{Adages} n° 677). Sur ces gâteaux de miel qui préservaient des dérangements intestinaux, voir aussi \italique{Thesm}. 570. On peut maintenant comprendre le bon mot de table qu'Athénée (\italique{Deipnosophistes}, I, 4d) rapporte : ἀμίδα δότω τις ἢ πλακοῦντά τις δότω, « Qu'on m'apporte un pot de chambre ou bien un gâteau ! » Voilà aussi pourquoi Aristophane fait intervenir dans ses \italique{Danaïdes} deux vendeurs opposés : celui de syrméa et celui d'ἐγκρίδες (Pollux, \italique{Onomasticon}, VII 199 = fr. 276(265) Kassel-Austin) : ἐν δ' Ἀριστοφάνης Δαναΐσι συρμαιοπῶλαι καὶ ἐγκριδοπῶλαι. Athénée (XIV, 645e) définit les encrides comme des gâteaux frits à l'huile et emmiellés : ΕΓΚΡΙΔΕΣ. πεμμάτιον ἑψόμενον ἐν ἐλαίῳ καὶ μετὰ τοῦτο μελιτούμενον.}.

Dans la \italique{Paix}, Trygée suggère au marchand d'armes menacé par la ruine de couper ses lances pour les proposer à la vente comme échalas. Autre idée : recycler ses casques d'une bien étrange manière (vers 1253-1254) :
\citComm{Τρυγαῖος – Πώλει βαδίζων αὐτὰ τοῖς Αἰγυπτίοις·\\
ἔστιν γὰρ ἐπιτήδεια συρμαίαν μετρεῖν.\\
Trygée – Va les vendre en Égypte.\\
Ils sont bons à mesurer la syrméa.}(C'est-à-dire à mesurer les \italique{effets} de la syrméa ; Trygée ne voit donc en ces casques que de vulgaires pots de chambre. En fait ils ne sont même bons à rien ; car Trygée fait entendre le proverbe κύματα μετρεῖν, compter les vagues, comble de l'inutilité\nc{Voir Érasme, \italique{Adages} (n° 345, \latin{Undas numeras}) ou Suétone, Περὶ βλασφημιῶν (éd. Taillardat), β 188.}.)

Même genre de détournement dans le présent passage des \italique{Thesmophories} : les Égyptiens se servent du Nil pour se \italique{nettoyer} (ici νοτίζειν = βαλανεύειν, σπογγίζειν) des désagréments de leur régime.

On peut conclure de ce passage qu'en grec le mot πέδον, le sol ; littéralement, ce sur quoi repose le pied, a fourni la même image qu'en français le \italique{fondement} et sert à désigner ἡ πυγή\nc{Noter que si πέδον, \italique{la plaine}, renvoie au verso ; πεδίον, \italique{la petite plaine}, elle, fait allusion au recto. Cf. \italique{Lysistrata}, v. 88 ; \italique{Oiseaux}, v. 507 avec le commentaire de Willems. Voir aussi l'emploi de λειμών au vers 171 du \italique{Cyclope} d'Euripide.}. 
Revenons à notre passage des \italique{Thesmophories} et proposons une traduction :\citComm{Μνησίλοχος – Νείλου μὲν αἵδε καλλιπάρθενοι ῥοαί,\\
ὃς ἀντὶ δίας ψακάδος Αἰγύπτου πέδον\\
λευκῆς νοτίζει μελανοσυρμαῖον λεών.\\
Mnésiloque – (Rôle d'Hélène.) « Du Nil ce sont ici les rives aux belles Nymphes,\\
Nil qui à la place de l'eau du ciel, en humecte le fondement\\
d'Égypte la blanche » et du peuple à la noire syrméa.}On comprend la réaction indignée de la gardienne de Mnésiloque (\italique{Thesm}. v. 858) :
\citComm{Κρίτυλλα – πανοῦργος εἶ νὴ τὴν Ἑκάτην τὴν φωσφόρον.\\
Critylla – Tu es un effronté, par Hécate la porte-torche.\\
(Elle lui porte un coup de sa torche.)}Cf. le vers 853 et le commentaire \italique{ad loc.} :
Critylla menaçait Mnésiloque d'un bon coup de flambeau. Et c'est au présent vers 858 qu'elle l'assène en évoquant Hécate \textbf{φωσ}φόρος.}
\comm{858}{τὴν Ἑκάτην τὴν φωσφόρον.}{Ce n'est pas par hasard que Critylla adjoint à Hécate l'épithète de φωσφόρος, \italique{porte-torche} : elle est elle-même équipée d'une telle torche et il est à parier qu'en jurant de la sorte, Critylla met sa menace du vers 853 à exécution en assénant à Mnésiloque le coup de flambeau promis. Cf. \italique{Lysistrata} vers 453 ; fr. 592.5 Kassel-Austin [$=$ Athénée XV, 701b] : ἐκφέρετε πεύκας κατ᾽ Ἀγάθωνα φωσφόρους.}
\comm{859–60}{}{\italique{Hélène}, v. 16-17a :\citComm{\difference{ἡμῖν} δὲ γῆ μὲν πατρὶς οὐκ ἀνώνυμος\\
Σπάρτη, πατὴρ δὲ Τυνδάρεως.}}
\comm{861}{Φρυνώνδας.}{Phrynondas et Eurybatos, deux coquins dont le nom était devenu appellatif, comme en français Cartouche et Mandrin. (Note de Willems ; qui a en tête le \italique{Protagoras} (327d) de Platon lorsque ces deux truands de la tradition populaire sont rapprochés.)\\
\italique{Souda}, φ 770 :\citComm{Φρυνώνδας· τῶν ἐπὶ πονηρίᾳ διαβεβοημένων. ὃς ξένος ὢν κατὰ τὰ Πελοποννησιακὰ διέτριβεν Ἀθήνησιν. Ἀριστοφάνης Ἀμφιάρεῳ· ὦ μιαρὲ καὶ Φρυνώνδα καὶ πονηρὲ σύ. ἐκ τούτου τοὺς πονηροὺς Φρυνώνδας καλοῦσι. κακοήθης ὁ Φρυνώνδας οὗτος καὶ πανοῦργος, μισθοῦ τινας μηχανὰς πραγμάτων ἐπὶ κακοῖς ῥᾳδίως συντιθείς, ὥς φησιν Εὔδημος. καὶ αὖθις· ἀλλ' οἶμαι οὔτε Φρυνώνδας οὔτ' Εὐρύβατος οὐδ' ἄλλος πώποτε τῶν πάλαι πονηρῶν, τοιοῦτος μάγος καὶ γόης.\\
\francais{Phrynondas : de ceux connus pour leur méchanceté. C'était un étranger qui vivait à Athènes au moment de la guerre du Péloponnèse. Aristophane dans son \italique{Amphiaraos} : « ô misérable et Phrynondas et vil que tu es ». On appelle les méchants de son nom : Phrynondas.
Ce Phrynondas [était] mauvais et méchant ; quelqu'un qui, pour de l'argent, élaborait facilement quelque intrigue pour de mauvaises fins, comme le dit Eudemos.
Et ailleurs : « mais je pense que ni Phrynondas ni Eurybatos ni aucun autre des vilains d'autrefois ne se sont jamais comportés comme un filou et un tricheur ».}}}
\comm{862}{}{\italique{Hélène}, v. 22 :\citComm{Ἑλένη δ᾽ ἐκλήθην. ἃ δὲ πεπόνθαμεν κακὰ}}
\comm{864–5}{}{\italique{Hélène}, v. 52-53a :\citComm{ψυχαὶ δὲ πολλαὶ δι᾽ ἔμ᾽ ἐπὶ Σκαμανδρίοις\\
ῥοαῖσιν ἔθανον· \difference{ἡ δὲ πάντα τλᾶσ᾽ ἐγὼ}}}
\comm{866–8}{}{\italique{Hélène}, v. 49 :\citComm{κἀγὼ μὲν ἐνθάδ᾽ εἴμ᾽, ὁ δ᾽ ἄθλιος πόσις}\italique{Hélène}, v. 56a : \citComm{τί δῆτ᾽ ἔτι ζῶ; }}
\comm{870}{μὴ ψεῦσον ὦ Ζεῦ τῆς ἐπιούσης ἐλπίδος.}{Voici la scholie au vers 870 :\citComm{σύνηθες τὸ σχῆμα. Μένανδρος·\\
\alinea οἶσθ᾽ ὃ ποίησον.\\
Σοφοκλῆς Πηλεῖ·\\
\alinea μὴ ψεῦσον ὦ Ζεῦ· μή μ᾽ ἕλης ἄνευ δορός.\\
\alinea οἶσθ᾽ οὖν ὃ δρᾶσον.\\
Εὐριπίδης Πολυΐδῳ.\\
\francais{La construction est usuelle. Ménandre : 
« Tu sais quoi faire. » 
Sophocle dans son \italique{Pelée} : 
« Ne me trompe pas, ô Zeus ; tu ne me prendras pas sans lutte. /
Tu sais donc ce qu'il te reste à faire. »
(Euripide dans son \italique{Polyidos}.)}}Cette scholie ne se peut comprendre qu'en la mettant en regard d'un passage du Περὶ διαλέκτων de Grégoire de Corinthe (\nieme{XII} siècle, p. 14 ed. Schaefer) ou de la scholie de l'\italique{Hécube} d'Euripide (vers 225 -- nos trois commentaires remontent assurément à la même source ; cf. aussi Suidas ψ 62). À propos de la tournure `οἶσθ᾽ οὖν ὃ δρᾶσον;', voici ce que le scholiaste de l'\italique{Hécube} écrit :\citComm{δρᾶσον: τὰ καλὰ τῶν ἀντιγράφων δρᾶσον ἔχει. οἱ Ἀττικοὶ χρῶνται τοῖς προστακτικοῖς ἀντὶ ὑποτακτικῶν ἀπαγορευτικῶν, καὶ αὐτοῖς τοῖς προστακτικοῖς ἀντὶ ὁριστικῶν.\\
τοῦ μὲν οὖν πρώτου τὸ Ἀριστοφάνους ἐν Θεσμοφοριαζούσαις (v. 870) παράδειγμα `μὴ ψεῦσον, ὦ Ζεῦ, τῆς ἐπιούσης ἡμέρας (sic).'\\
τοῦ δευτέρου δὲ Μένανδρος `οἶσθ' ὅ τι ποίησον' ἀντὶ 
τοῦ `ὅ τι ποιήσεις'.\\
τὸν αὐτὸν τρόπον καὶ Εὐριπίδης `οἶσθ᾿ οὖν ὃ δρᾶσον' ἀντὶ τοῦ `δράσεις' ἐξήνεγκεν.\\
δρᾶσον: les meilleures copies ont `δρᾶσον'.
Les Attiques recourent aux impératifs à la place de subjonctifs de prohibition. Et de nouveau aux (je lis : καὶ αὖ τοῖς) impératifs à la place d'indicatifs.\\
L'exemple d'Aristophane sur le premier point : `μὴ ψεῦσον κτλ.'\\
Ménandre sur le second point : `οἶσθ' ὅ τι ποίησον' à la place de `ὅ τι ποιήσεις'.\\
Euripide a le même tour : `οἶσθ᾿ οὖν ὃ δρᾶσον' à la place de `δράσεις'.}Noter que ce tour d'Euripide se trouve justement dans l'\italique{Hélène} (vers 325) :\citComm{Χορός --- οἶσθ᾽ οὖν ὃ δρᾶσον; μνήματος λιποῦσ᾽ ἕδραν.\\
\francais{Hé bien, sais-tu ce que tu dois faire ? Quitte ton poste près du tombeau.}}}
\comm{871}{}{\italique{Hélène}, v. 68 :\citComm{Τεῦκρος --- τίς τῶνδ᾽ ἐρυμνῶν δωμάτων ἔχει κράτος;}\italique{Hélène}, v. 449 :\citComm{Μενελέως --- ναυαγὸς \difference{ἥκω} ξένος, \difference{ἀσύλητον γένος}.}}
\comm{874–8}{}{\italique{Hélène}, v. 460 : \citComm{Πρωτεὺς τάδ᾽ οἰκεῖ δώματ᾽, Αἴγυπτος δὲ γῆ.}}
\comm{874–6}{ποίου Πρωτέως; ... ἐπεὶ τέθνηκε Πρωτέας ἔτη δέκα.}{Elle confond Protée avec Protéas\nc{Identiques au génitif.}, stratège athénien dont il est question dans Thucydide\nc{Thucydide I, 45.2 ; II, 23.2.}. (Note de Willems.)}
